
\section{Metodologia}
A metodologia adotada neste trabalho envolve a análise comparativa de diferentes linguagens de programação e bibliotecas utilizadas em visão computacional. O processo metodológico foi dividido nas seguintes etapas:

\subsection{Seleção das linguagens e bibliotecas}

Foram selecionadas as principais linguagens de programação (Java, Python, C, C++) e bibliotecas (OpenCV, TensorFlow, YOLO, JavaCV, VisionWorks, BoofCV, OpenIMAJ, Halcon, SimpleCV) amplamente utilizadas em visão computacional.

\subsection{Definição dos critérios de avaliação}

Foram estabelecidos critérios para a avaliação das linguagens e bibliotecas, incluindo facilidade de uso, curva de aprendizado, comunidade de suporte, desempenho, compatibilidade com diferentes plataformas e recursos disponíveis para visão computacional.

Este critério avalia a complexidade sintática e o esforço cognitivo necessários para a implementação das soluções. A análise baseou‑se em duas métricas principais:

Métrica Quantitativa (LOC): Utilizou‑se a contagem de Linhas de Código (Lines of Code - LOC) necessárias para implementar os experimentos em cada linguagem.

Métrica Qualitativa (Escala Likert): Foi estabelecida uma escala de 1 a 5 (onde 1 = Insuficiente e 5 = Excelente) para classificar três aspectos: (i) clareza da documentação oficial; (ii) disponibilidade de tutoriais/comunidade; e (iii) facilidade de configuração do ambiente (setup).

\subsubsection{Desempenho}
O desempenho tem o objetivo de medir a eficiência das linguagens e bibliotecas na execução de tarefas comuns em visão computacional, sob cenários controlados. Foram realizados testes práticos para avaliar o tempo de execução, uso de memória e capacidade de lidar com grandes volumes de dados.

\subsubsection{Compatibilidade com Diferentes Plataformas}
A compatibilidade com diferentes plataformas tem o objetivo de verificar a capacidade das linguagens e bibliotecas de serem executadas em diversos sistemas operacionais (Windows, Linux, macOS) e dispositivos (desktops, dispositivos móveis, sistemas embarcados). Foram analisadas as dependências necessárias e a facilidade de instalação em cada ambiente.

\subsubsection{Recursos Disponíveis para Visão Computacional}
O critério de recursos disponíveis para visão computacional tem o objetivo de avaliar a variedade e qualidade das funcionalidades oferecidas por cada biblioteca, como algoritmos de detecção   reconhecimento de objetos, processamento de imagens, aprendizado de máquina, entre outros. Foram analisadas as capacidades específicas de cada biblioteca em relação às necessidades comuns em projetos de visão computacional.

\subsection{Coleta de dados}

Para garantir a reprodutibilidade dos testes, foram adotadas múltiplas fontes de consulta, incluindo a documentação oficial das bibliotecas, repositórios de código aberto e artigos técnicos especializados. Nos experimentos práticos, definiram-se dois conjuntos de tarefas-padrão para isolar o desempenho da linguagem:

\begin{enumerate}
    \item \textbf{Detecção de Mãos utilizando Classificadores Haar Cascade:} Esta tarefa serve como \textit{baseline} de desempenho, utilizando o classificador em cascata disponível na biblioteca OpenCV. Para garantir a paridade entre as linguagens (Java, Python e C++), foram fixados rigidamente os mesmos hiperparâmetros de detecção:
\begin{itemize}
    \item \textbf{Modelo:} Arquivo XML pré-treinado para detecção de mãos (\textit{hand.xml});
    \item \textbf{Parâmetros de Escala:} \textit{scaleFactor} de 1.1 (redução de 10\% a cada etapa da pirâmide de imagem) e \textit{minNeighbors} igual a 1 (para maximizar a sensibilidade, ainda que com risco de falsos positivos);
    \item \textbf{Tamanho Mínimo:} Janela de detecção mínima fixada em $150 \times 150$ pixels para ignorar ruídos menores;
    \item \textbf{Saída:} Desenho do retângulo delimitador (\textit{Bounding Box}) sobre a imagem original e persistência do resultado em disco.
\end{itemize}

    \item \textbf{Detecção de Mãos baseada em Processamento de Imagem (Baixo Nível):} Esta tarefa avalia a eficiência da linguagem na manipulação direta de matrizes e cálculos geométricos, sem o uso de classificadores de "caixa preta". O algoritmo foi implementado manualmente seguindo a lógica de segmentação e geometria:
    \begin{itemize}
        \item \textbf{Segmentação da Pele:} Conversão do espaço de cor RGB para \textbf{YCrCb}, seguida de limiarização (\textit{Thresholding}) com valores pré-definidos para isolar os pixels de pele;
        \item \textbf{Refinamento da Máscara:} Aplicação de Filtro Gaussiano, seguido de operações morfológicas de Abertura e Fechamento (\textit{Morphological Open/Close}) e \textit{Median Blur} para eliminação de ruído;
        \item \textbf{Análise Geométrica:} Extração do maior contorno da cena e cálculo do Fecho Convexo (\textit{Convex Hull});
        \item \textbf{Classificação de Gestos:} Identificação dos Defeitos de Convexidade (\textit{Convexity Defects}) e aplicação da Lei dos Cossenos para filtrar ângulos entre os dedos. A contagem de dedos determina a classe do gesto (ex: "Palm", "Fist", "Peace").
    \end{itemize}
\end{enumerate}

O ambiente de teste utilizado foi o seguinte:


\begin{table}[h]
\centering
\caption{Ambiente de Teste e Versões}
\begin{tabular}{l l}
\hline
\textbf{Linguagem/Biblioteca} & \textbf{Versão} \\
\hline
Python & 3.13.5 \\
Java & 20 \\
C++ (GCC) & 12.2.0 \\
\hline
\multicolumn{2}{l}{\textbf{Hardware - Windows}} \\
\hline
Processador & AMD Ryzen 5 3600X \\
Memória RAM & 32 GB \\
GPU & NVIDIA GeForce RTX 2060 \\
Sistema Operacional & Windows 10 Pro 64 bits \\
\hline
\end{tabular}
\end{table}


Linux: 

MacOS: 

Cada experimento foi repetido várias vezes, e os resultados foram agregados por média e desvio‑padrão para reduzir a variabilidade. Todo o código de teste, scripts de medição e procedimentos de pré‑processamento foram versionados e disponibilizados no repositório do projeto para garantir reprodutibilidade. Finalmente, registraram‑se limitações e potenciais fontes de viés (diferenças de implementação, otimizações específicas de hardware, dependências) e considerações éticas relativas ao uso dos conjuntos de dados .

Conjuntos de Dados Utilizados:
\begin{itemize}
    \item Credits for \href{https://www.kaggle.com/datasets/ritikagiridhar/2000-hand-gestures?resource=download}{Hand gestures dataset} from Ritika Giridhar at Kaggle.com
\end{itemize}

\subsection{Análise comparativa}
Os dados coletados foram analisados comparativamente, utilizando técnicas estatísticas e gráficos para identificar as vantagens e desvantagens de cada linguagem e biblioteca em relação aos critérios estabelecidos. A análise incluiu a criação de tabelas comparativas, gráficos de desempenho e discussões qualitativas sobre as características observadas.



